\input{header}

\title{Data Cleaning \& Claude Code for Research}
\subtitle{ECON 692 -- Applied Economics Seminar}
\author{Andrew Hobbs}
\institute{University of San Francisco}
\date{February 26, 2026}

\begin{document}

\begin{frame}[plain]
  \titlepage
\end{frame}

% ══════════════════════════════════════════════════════════
%  Agenda & Check-in
% ══════════════════════════════════════════════════════════

\begin{frame}{Today's Agenda}
  \begin{tabular}{@{}r@{\hspace{8pt}}l@{\hspace{16pt}}l@{}}
    \toprule
    \textbf{Time} & \textbf{Duration} & \textbf{Activity} \\
    \midrule
    4:35 & 5 min  & Welcome \& check-in \\
    4:40 & 10 min & Claude Code setup verification \\
    4:50 & 40 min & Claude Code intro \& CLAUDE.md workshop \\
    5:30 & 10 min & Break \\
    5:40 & 30 min & Data cleaning best practices \\
    6:10 & 10 min & Break \\
    6:20 & 45 min & Claude Code for data tasks \\
    7:05 & 35 min & Work sprint: clean your own data \\
    7:40 & 15 min & Share-outs \\
    7:55 & 5 min  & Wrap-up \& next week \\
    \bottomrule
  \end{tabular}
\end{frame}

\begin{frame}{Announcement: Next Week's Schedule}
  \textbf{No regular class on Thursday, March 5.} Instead:

  \vspace{8pt}

  \textbf{1. Morning session: Thursday, March 5, 10:00am--12:00pm}
  \begin{itemize}
    \item Replacement class covering Week 6 material (EDA and early figures)
    \item Location: \textbf{Lone Mountain 244B}
  \end{itemize}

  \vspace{8pt}

  \textbf{2. Evening seminar: Thursday, March 5, 6:30--8:30pm at Stanford}
  \begin{itemize}
    \item \textbf{Bay Area Tech Economics Seminar} with Andrey Fradkin (BU / Amazon)
    \item Topic: economics of AI -- LLM market competition and the ``Coasean Singularity''
    \item Location: Simonyi Conference Center, 389 Jane Stanford Way
  \end{itemize}

  \vspace{8pt}

  Both are \textbf{strongly encouraged}. I won't take attendance for either.
\end{frame}

\begin{frame}{Check-in}
  \begin{itemize}
    \item How is your data looking? Any surprises when you loaded it?
    \item Has anyone started cleaning or merging datasets?
    \item Did everyone get the Claude Code install instructions?
  \end{itemize}
  \vspace{12pt}
  \begin{block}{Reminder}
    No major milestone this week -- use this time to invest in your tools and data pipeline.
    \textbf{Empirical Strategy Draft} due \textbf{Friday, March 6}.
  \end{block}
\end{frame}

\begin{frame}{Learning Objectives}
  Last week you identified \textit{what} to estimate and drew your DAG. This week you set up your tools, clean your data, and push toward \textbf{initial results}.

  \vspace{6pt}

  By the end of today, you will be able to:

  \vspace{4pt}

  \begin{enumerate}
    \item Set up and use \textbf{Claude Code} as a research assistant in your terminal
    \item Write a \textbf{CLAUDE.md} file that gives Claude Code context about your project
    \item Apply \textbf{tidy data principles} and build a reproducible data pipeline
    \item Use Claude Code to \textbf{clean, model, and visualize} -- aiming for initial results
  \end{enumerate}
\end{frame}

% ══════════════════════════════════════════════════════════
%  SECTION 1: Claude Code
% ══════════════════════════════════════════════════════════

\section{Claude Code}

\begin{frame}[fragile]{What Is Claude Code?}
  \begin{columns}[T]
    \begin{column}{0.47\textwidth}
      Claude Code is an AI assistant that lives in your terminal.

      \vspace{8pt}

      \begin{itemize}
        \item Reads your entire project -- files, folders, git history
        \item Writes and edits code for you
        \item Runs commands and checks output
        \item Remembers your conventions via \texttt{CLAUDE.md}
      \end{itemize}

      \vspace{8pt}

      Think of it as a collaborator who can read your whole project and write code on the spot.
    \end{column}
    \begin{column}{0.50\textwidth}
\begin{lstlisting}[style=terminal, basicstyle=\ttfamily\footnotesize\color{darkgray}]
$ claude

> Read my dataset and tell me
  what needs cleaning.

Claude: I found 3 issues in
  data/raw/survey.csv:
  1. 847 missing values in income
  2. State codes are inconsistent
     (CA, Calif., California)
  3. 23 duplicate rows
\end{lstlisting}
    \end{column}
  \end{columns}
\end{frame}

\begin{frame}{Why Claude Code for Research?}
  \textbf{The research workflow problem:}

  \begin{itemize}
    \item You spend 60--80\% of research time on data cleaning, not analysis
    \item Repetitive tasks: recoding, reshaping, merging, validating
    \item Context switching: Stack Overflow $\to$ code $\to$ run $\to$ debug $\to$ repeat
  \end{itemize}

  \vspace{10pt}

  \textbf{What Claude Code changes:}

  \begin{itemize}
    \item Describe what you want in plain English
    \item Claude Code reads your data, writes the code, runs it
    \item You review the output and iterate
    \item Every step is in a script -- fully reproducible
  \end{itemize}

  \vspace{8pt}

  \textit{This is not about replacing your thinking. It is about automating the tedious parts so you can focus on your research design.}
\end{frame}

\begin{frame}[fragile]{Setup Verification}
\begin{lstlisting}[style=terminal]
# Check installation
claude --version

# First launch (log in with your Pro account)
claude

# Test it works
> What directory am I in?
\end{lstlisting}

  \vspace{4pt}

  \small
  \textbf{Troubleshooting:} \texttt{command not found} $\to$ re-run the installer from \texttt{docs.anthropic.com} \quad|\quad
  Login issues $\to$ check your email for the Pro invite \quad|\quad Windows $\to$ use WSL

  \vspace{4pt}

  \normalsize
  \textbf{\textcolor{usfgreen}{Take 5 minutes now.}} Open your terminal and verify Claude Code runs. Raise your hand if you need help.
\end{frame}

\begin{frame}{How Claude Code Sees Your Project}
  \centering
  \begin{tikzpicture}[>=stealth, scale=0.8, font=\small,
      folder/.style={draw, rounded corners, fill=codebg, minimum width=2cm, minimum height=0.6cm},
      leaf/.style={draw, rounded corners, fill=codebg, minimum width=1.6cm, minimum height=0.55cm, font=\footnotesize},
      special/.style={draw, rounded corners, fill=usfgreen!15, draw=usfgreen, thick, minimum width=2cm, minimum height=0.6cm}]

    % Project root
    \node[folder] (root) at (0, 3) {\texttt{my-capstone/}};

    % Children
    \node[special] (claude) at (-5, 1.5) {\texttt{CLAUDE.md}};
    \node[folder] (data) at (-1.8, 1.5) {\texttt{data/}};
    \node[folder] (code) at (1.8, 1.5) {\texttt{code/}};
    \node[folder] (output) at (5, 1.5) {\texttt{output/}};

    % Sub-children (wider spacing, smaller boxes)
    \node[leaf] (raw) at (-3, 0) {\texttt{raw/}};
    \node[leaf] (clean) at (-0.6, 0) {\texttt{cleaned/}};
    \node[leaf] (scripts) at (1.8, 0) {\texttt{*.R}};
    \node[leaf] (figs) at (5, 0) {\texttt{figures/}};

    % Edges
    \draw[->] (root) -- (claude);
    \draw[->] (root) -- (data);
    \draw[->] (root) -- (code);
    \draw[->] (root) -- (output);
    \draw[->] (data) -- (raw);
    \draw[->] (data) -- (clean);
    \draw[->] (code) -- (scripts);
    \draw[->] (output) -- (figs);

    % Annotation
    \node[usfgreen, font=\footnotesize, align=center] at (0, -1.1) {
      Claude Code reads \textbf{all} of these files\\
      and understands how they relate to each other.\\[4pt]
      \textbf{CLAUDE.md} is the first thing it reads every session.
    };
  \end{tikzpicture}
\end{frame}

% ══════════════════════════════════════════════════════════
%  SECTION 2: CLAUDE.md & Project Setup
% ══════════════════════════════════════════════════════════

\section{CLAUDE.md \& Project Setup}

\begin{frame}{What Is CLAUDE.md?}
  A \textbf{CLAUDE.md} file is a plain-text instruction manual for Claude Code.

  \vspace{8pt}

  It lives in the root of your project and tells Claude Code:

  \begin{itemize}
    \item What your project is about
    \item How your folders are organized
    \item What commands to run (R, Python, LaTeX, etc.)
    \item Conventions and rules to follow
    \item What data sources you are using
  \end{itemize}

  \vspace{8pt}

  Without CLAUDE.md, Claude Code starts every session from scratch. With it, Claude Code understands your project \textbf{immediately}.
\end{frame}

\begin{frame}[fragile]{Anatomy of a Research CLAUDE.md}
\begin{lstlisting}[style=terminal, basicstyle=\ttfamily\scriptsize\color{darkgray}]
# CLAUDE.MD -- [Your Name] Capstone Project

## Project
- Topic: Effect of [X] on [Y] using [method]
- Language: R (tidyverse, fixest)
- Data: [source], [unit of observation], [time period]

## Folder Structure
- data/raw/     -- original data (never modify)
- data/cleaned/ -- processed data
- code/         -- R scripts (01_clean.R, 02_merge.R, ...)
- output/       -- figures and tables

## Commands
Rscript code/01_clean.R     # cleaning
Rscript code/03_analysis.R  # estimation

## Conventions
- snake_case variable names; NA for missing (never 999)
\end{lstlisting}
\end{frame}

\begin{frame}{CLAUDE.md Workshop}
  \textbf{\textcolor{usfgreen}{15 minutes.}} Create a CLAUDE.md for your capstone project.

  \vspace{6pt}

  \begin{enumerate}
    \item Open your capstone repository in your terminal
    \item Create a file called \texttt{CLAUDE.md} in the project root
    \item Fill in these sections (use the template from the previous slide):
    \begin{itemize}
      \item \textbf{Project}: one-sentence description, language, data source
      \item \textbf{Folder structure}: your actual directories
      \item \textbf{Commands}: how to run your main scripts
      \item \textbf{Conventions}: any rules you want Claude Code to follow
    \end{itemize}
    \item Launch Claude Code in your project: \texttt{cd your-repo \&\& claude}
    \item Ask: ``Read my CLAUDE.md and summarize what you know about my project.''
  \end{enumerate}

  \vspace{4pt}

  \textit{I will walk around to help. This file will save you hours over the semester.}
\end{frame}

\begin{frame}[fragile]{Rules Files}
  For conventions too detailed for CLAUDE.md, use \textbf{rules files} in \texttt{.claude/rules/}:

  \vspace{4pt}

  \textbf{Example:} \texttt{.claude/rules/data-conventions.md}

\begin{lstlisting}[style=terminal, basicstyle=\ttfamily\footnotesize\color{darkgray}]
# Data Conventions
- Raw data files are NEVER modified
- All cleaning steps are scripted in code/
- Missing values: NA only (never 999, -1, or "")
- State identifiers: always 2-digit FIPS codes
- Date format: YYYY-MM-DD
\end{lstlisting}

  \vspace{4pt}

  Claude Code reads rules files automatically every session. They persist your standards across conversations.

  \vspace{4pt}

  \textit{Think of rules files as coding standards that Claude Code enforces for you.}
\end{frame}

\begin{frame}{Giving Good Instructions}
  Claude Code is powerful, but \textbf{garbage in, garbage out} still applies.

  \vspace{6pt}

  \begin{columns}[T]
    \begin{column}{0.47\textwidth}
      \textbf{\textcolor{red}{Vague (avoid):}}
      \begin{itemize}
        \item ``Clean my data''
        \item ``Fix the problems''
        \item ``Make it work''
      \end{itemize}
    \end{column}
    \begin{column}{0.47\textwidth}
      \textbf{\textcolor{usfgreen}{Specific (better):}}
      \begin{itemize}
        \item ``Read data/raw/cps.csv. Drop rows where income $<$ 0. Recode state names to FIPS codes.''
        \item ``Merge enrollment.csv and outcomes.csv on student\_id. Report unmatched rows.''
      \end{itemize}
    \end{column}
  \end{columns}

  \vspace{6pt}

  \textbf{Three principles:}
  \begin{enumerate}
    \item \textbf{Be specific} about files, variables, and expected output
    \item \textbf{State your goal}, not just the steps
    \item \textbf{Verify}: ask Claude Code to explain what it did, then check
  \end{enumerate}
\end{frame}

% ══════════════════════════════════════════════════════════
%  SECTION 3: Data Cleaning Principles
% ══════════════════════════════════════════════════════════

\section{Data Cleaning Principles}

\begin{frame}{The Data Pipeline}
  \centering
  \begin{tikzpicture}[>=stealth, scale=0.75, font=\small,
      stage/.style={draw, rounded corners, minimum width=2.4cm, minimum height=1cm, align=center, thick}]

    \node[stage, fill=red!10, draw=red!40] (raw) at (0, 0)
      {Raw Data\\{\footnotesize (never touch)}};
    \node[stage, fill=usfgold!15, draw=usfgold] (clean) at (6, 0)
      {Cleaned Data\\{\footnotesize (scripted)}};
    \node[stage, fill=usfgreen!15, draw=usfgreen] (analysis) at (12, 0)
      {Analysis\\{\footnotesize (reproducible)}};

    \draw[->, thick, gray] (raw) -- node[above, font=\footnotesize] {\texttt{01\_clean.R}} (clean);
    \draw[->, thick, gray] (clean) -- node[above, font=\footnotesize] {\texttt{02\_merge.R \ldots}} (analysis);

    \node[font=\footnotesize, align=center, gray] at (6, -1.5) {
      Every transformation is in a script. Delete \texttt{cleaned/} and re-run from scratch.
    };
  \end{tikzpicture}

  \vspace{8pt}

  \begin{block}{Golden Rule of Reproducibility}
    Raw data is sacred. Scripts are the single source of truth. If you cannot
    regenerate your cleaned data from raw data + scripts alone, your pipeline is broken.
  \end{block}

  \vspace{2pt}

  \textit{Could someone reproduce your current results from raw data + scripts alone? Where does the chain break?}
\end{frame}

\begin{frame}{Tidy Data Principles}
  \small
  \textbf{Tidy data} (Wickham, 2014): every \textbf{variable} is a column, every \textbf{observation} is a row, every \textbf{observational unit} is a table.

  \vspace{6pt}

  \begin{columns}[T]
    \begin{column}{0.47\textwidth}
      \textbf{\textcolor{red}{Messy}} (wide):

      \vspace{4pt}

      \footnotesize
      \begin{tabular}{@{}lccc@{}}
        \toprule
        State & GDP\_19 & GDP\_20 & GDP\_21 \\
        \midrule
        CA & 3.1T & 3.0T & 3.4T \\
        TX & 1.8T & 1.7T & 1.9T \\
        \bottomrule
      \end{tabular}
    \end{column}
    \begin{column}{0.47\textwidth}
      \textbf{\textcolor{usfgreen}{Tidy}} (long):

      \vspace{4pt}

      \footnotesize
      \begin{tabular}{@{}lcc@{}}
        \toprule
        State & Year & GDP \\
        \midrule
        CA & 2019 & 3.1T \\
        CA & 2020 & 3.0T \\
        CA & 2021 & 3.4T \\
        TX & 2019 & 1.8T \\
        \bottomrule
      \end{tabular}
    \end{column}
  \end{columns}

  \vspace{6pt}

  \normalsize
  Most panel data methods (DiD, TWFE, event studies) require tidy (long) format.

  \vspace{2pt}

  R: \texttt{pivot\_longer()} / \texttt{pivot\_wider()} \quad|\quad Python: \texttt{pd.melt()} / \texttt{pd.pivot()}
\end{frame}

\begin{frame}[fragile]{Common Data Problems}
  \small
  \begin{tabular}{@{}l>{\raggedright\arraybackslash}p{4.5cm}>{\raggedright\arraybackslash}p{4.2cm}@{}}
    \toprule
    \textbf{Problem} & \textbf{Example} & \textbf{R Fix} \\
    \midrule
    Missing values &
      Income coded as 999, $-$1, or blank &
      \texttt{mutate(inc = na\_if(inc, 999))} \\[4pt]
    Duplicates &
      Same observation appears twice &
      \texttt{distinct()} or \texttt{filter(n() > 1)} \\[4pt]
    Inconsistent coding &
      ``California'', ``CA'', ``calif.'' &
      \texttt{case\_when()} or lookup join \\[4pt]
    Type mismatch &
      Numeric column read as character &
      \texttt{mutate(x = as.numeric(x))} \\[4pt]
    Wide format &
      Years as column names &
      \texttt{pivot\_longer()} \\
    \bottomrule
  \end{tabular}

  \vspace{4pt}

  \footnotesize
  Python: \texttt{df.replace(999, np.nan)}, \texttt{df.drop\_duplicates()}, \texttt{pd.melt()}, \texttt{df.astype()}

  \vspace{2pt}

  \normalsize
  \textbf{\textcolor{usfgreen}{Warning:}} Claude Code may handle these differently than you expect. Be explicit and check what it did.
\end{frame}

\begin{frame}[fragile]{Merging Datasets}
  \begin{columns}[T]
    \begin{column}{0.44\textwidth}
      \centering
      \vspace{4pt}
      \begin{tikzpicture}[scale=0.55]
        % Left join
        \node[font=\footnotesize\bfseries] at (2, 3.7) {\texttt{left\_join()}};
        \fill[usfgreen!20] (1.2, 2) circle (1);
        \draw[thick] (1.2, 2) circle (1);
        \draw[thick] (2.8, 2) circle (1);
        \node[font=\scriptsize] at (0.7, 2) {A};
        \node[font=\scriptsize] at (3.3, 2) {B};
        \node[font=\scriptsize, gray] at (2, 0.5) {Keep all of A};

        % Inner join
        \node[font=\footnotesize\bfseries] at (7, 3.7) {\texttt{inner\_join()}};
        \draw[thick] (6.2, 2) circle (1);
        \draw[thick] (7.8, 2) circle (1);
        \begin{scope}
          \clip (6.2, 2) circle (1);
          \fill[usfgreen!20] (7.8, 2) circle (1);
        \end{scope}
        \draw[thick] (6.2, 2) circle (1);
        \draw[thick] (7.8, 2) circle (1);
        \node[font=\scriptsize] at (5.7, 2) {A};
        \node[font=\scriptsize] at (8.3, 2) {B};
        \node[font=\scriptsize, gray] at (7, 0.5) {Keep only matches};
      \end{tikzpicture}
    \end{column}
    \begin{column}{0.53\textwidth}
      \textbf{Key decisions:}
      \begin{itemize}
        \item What is the \textbf{key}? (state + year? student\_id?)
        \item Are keys \textbf{unique}? (many-to-many is almost always a bug)
        \item What happens to \textbf{non-matches}?
      \end{itemize}

      \vspace{4pt}

      \textbf{Always check your merge:}

\begin{lstlisting}[style=terminal, basicstyle=\ttfamily\footnotesize\color{darkgray}]
# Rows that failed to match
anti_join(A, B, by = "state_fips")
\end{lstlisting}

      \vspace{2pt}

      \footnotesize
      Python: \texttt{pd.merge(A, B, how="left", indicator=True)}
    \end{column}
  \end{columns}
\end{frame}

\begin{frame}[fragile]{Data Validation}
  \small
  After cleaning, \textbf{verify your data} before analysis:

  \vspace{4pt}

  \begin{columns}[T]
    \begin{column}{0.47\textwidth}
      \begin{enumerate}
        \item \textbf{Row counts} -- expected $n$?
        \item \textbf{Key uniqueness} -- no duplicate (unit, time)?
        \item \textbf{Range checks} -- plausible values?
        \item \textbf{Missing patterns} -- random or systematic?
        \item \textbf{Balance} -- treatment \& control in same periods?
      \end{enumerate}
    \end{column}
    \begin{column}{0.50\textwidth}
\begin{lstlisting}[style=terminal, basicstyle=\ttfamily\scriptsize\color{darkgray}]
# Quick validation in R
nrow(df)
df |> count(state, year) |>
  filter(n > 1)
summary(df$income)
df |> group_by(treated) |>
  summarize(n = n(),
    pct_na = mean(is.na(income)))
\end{lstlisting}
    \end{column}
  \end{columns}
\end{frame}

\begin{frame}{File Organization Best Practices}
  \begin{columns}[T]
    \begin{column}{0.47\textwidth}
      \textbf{Recommended structure:}

      \vspace{4pt}

      \small
      \begin{tabular}{@{}ll@{}}
        \texttt{data/raw/}     & Original files \\
        \texttt{data/cleaned/} & Processed data \\
        \texttt{code/01\_clean.R}  & Cleaning \\
        \texttt{code/02\_merge.R}  & Merging \\
        \texttt{code/03\_analysis.R} & Estimation \\
        \texttt{output/}       & Figures \& tables \\
        \texttt{paper/}        & Written report \\
      \end{tabular}
    \end{column}
    \begin{column}{0.47\textwidth}
      \textbf{Why numbered scripts?}
      \begin{itemize}
        \item Clear execution order
        \item Anyone can reproduce your results by running 01, 02, 03 in sequence
        \item Claude Code understands the pipeline when it reads your CLAUDE.md
      \end{itemize}
    \end{column}
  \end{columns}

  \vspace{8pt}

  \begin{block}{Reminder}
    Add \texttt{data/raw/} to \texttt{.gitignore} if files are large.
    Keep scripts in Git. Never commit cleaned data -- regenerate it from scripts.
  \end{block}
\end{frame}

% ══════════════════════════════════════════════════════════
%  SECTION 4: Claude Code for Data Tasks
% ══════════════════════════════════════════════════════════

\section{Claude Code for Data Tasks}

\begin{frame}{Live Demo}
  \textbf{I will now demonstrate a complete data cleaning workflow using Claude Code.}

  \vspace{10pt}

  \begin{enumerate}
    \item \textbf{Load \& inspect}: ``Read the data and describe what you see.''
    \item \textbf{Diagnose}: ``What needs to be cleaned before I can run my model?''
    \item \textbf{Clean}: ``Fix codes, handle missing values, reshape to long.''
    \item \textbf{Merge}: ``Merge these two datasets. Report unmatched rows.''
    \item \textbf{Validate}: ``Check for duplicates, summarize the panel structure.''
    \item \textbf{Save}: ``Write the cleaned data and commit the script to Git.''
  \end{enumerate}

  \vspace{10pt}

  \textbf{Key patterns to watch for:}
  \begin{itemize}
    \item Describe and diagnose \textit{before} asking Claude Code to fix
    \item Always ask for a \textit{script}, not ad hoc code -- scripts are reproducible
    \item Validate after every step -- Claude Code catches things you might miss
  \end{itemize}
\end{frame}

\begin{frame}{Step 1: Load \& Inspect}
  \textbf{Before touching anything, understand what you have.}

  \vspace{8pt}

  Ask Claude Code to read your raw data and report:
  \begin{itemize}
    \item How many rows and columns?
    \item What are the variable names and types?
    \item How many missing values per column?
    \item Are there obvious data quality issues?
  \end{itemize}

  \vspace{8pt}

  \textbf{\textcolor{usfgreen}{Key pattern:}} always \textit{describe and diagnose} before asking Claude Code to fix anything. If it does not know what the data looks like, it will guess -- and guess wrong.

  \vspace{8pt}

  \textit{This is also your chance to verify Claude Code read the right file and understands its structure.}
\end{frame}

\begin{frame}{Step 2: Clean \& Reshape}
  \textbf{Turn diagnoses into a scripted pipeline.}

  \vspace{8pt}

  Tell Claude Code \textit{exactly} what to do -- and ask it to write a script, not run ad hoc commands:
  \begin{itemize}
    \item Which variables need recoding, and \textit{how}?
    \item How should missing values be handled? (drop? recode? flag?)
    \item Does the data need reshaping? (wide $\to$ long for panel methods)
    \item Where should the output be saved?
  \end{itemize}

  \vspace{8pt}

  \textbf{\textcolor{usfgreen}{Key pattern:}} ``Write a script called \texttt{code/01\_clean.R} that does X, Y, Z and saves to \texttt{data/cleaned/}.'' Scripts are reproducible; ad hoc commands are not.
\end{frame}

\begin{frame}{Step 3: Merge \& Validate}
  \textbf{Combine datasets and immediately check the result.}

  \vspace{8pt}

  When merging, always specify:
  \begin{itemize}
    \item The \textbf{key} variables to join on
    \item The \textbf{join type} (left, inner, full) -- do not let Claude Code choose
    \item What to do with \textbf{non-matches} -- report them with \texttt{anti\_join()}
  \end{itemize}

  \vspace{8pt}

  After every merge, validate:
  \begin{itemize}
    \item Did the row count change as expected?
    \item Are there duplicate (unit, time) combinations?
    \item Are key variables complete across treatment and control groups?
  \end{itemize}

  \vspace{8pt}

  \textbf{\textcolor{usfgreen}{Key pattern:}} Claude Code sometimes catches issues you would miss. Always ask it to validate -- but verify its conclusions too.
\end{frame}

\begin{frame}{Step 4: Save \& Commit}
  \textbf{Lock in your work with Git.}

  \vspace{8pt}

  Once your pipeline runs cleanly:
  \begin{itemize}
    \item Commit your \textbf{scripts} to Git (not the cleaned data -- regenerate it)
    \item Add \texttt{data/cleaned/} to \texttt{.gitignore} if files are large
    \item Update your \textbf{CLAUDE.md} to describe the pipeline
    \item Write a clear commit message: what changed and why
  \end{itemize}

  \vspace{8pt}

  \textbf{Test reproducibility:} delete \texttt{data/cleaned/}, re-run your scripts, and confirm you get the same result. If you can't, the pipeline is broken.

  \vspace{8pt}

  \textit{You can ask Claude Code to do all of this: ``Commit my cleaning scripts, add cleaned data to .gitignore, and update CLAUDE.md with the pipeline description.''}
\end{frame}

\begin{frame}{Hands-On: Clean Your Data with Claude Code}
  \textbf{\textcolor{usfgreen}{20 minutes.}} Use Claude Code on your own data.

  \vspace{4pt}

  \textbf{Setup:}
  \begin{enumerate}
    \item Navigate to your capstone project: \texttt{cd \textasciitilde/path/to/your-repo}
    \item Launch Claude Code: \texttt{claude}
  \end{enumerate}

  \vspace{4pt}

  \textbf{Tasks (in order):}
  \begin{enumerate}
    \item Ask Claude Code to read one of your raw data files and describe it
    \item Ask it to identify data quality issues
    \item Ask it to write a cleaning script (\texttt{code/01\_clean.R} or \texttt{.py})
    \item Run the script and verify the output
    \item Ask Claude Code to validate the cleaned data
  \end{enumerate}

  \vspace{4pt}

  \textbf{If you do not have data yet:} ask Claude Code to help you find and download a public dataset relevant to your topic.

  \vspace{4pt}

  \textit{Raise your hand if you get stuck.}
\end{frame}

\begin{frame}{Common Claude Code Patterns for Research}
  \footnotesize
  \begin{tabular}{@{}l>{\raggedright\arraybackslash}p{8cm}@{}}
    \toprule
    \textbf{Task} & \textbf{What to say} \\
    \midrule
    Describe data &
      ``Read data/raw/file.csv and summarize: rows, columns, types, missing values.'' \\[2pt]
    Clean a variable &
      ``In 01\_clean.R, recode state names to 2-letter abbreviations using a lookup table.'' \\[2pt]
    Reshape &
      ``Pivot the wide GDP columns to long format with year and gdp columns.'' \\[2pt]
    Merge &
      ``Left join enrollment.csv and outcomes.csv on student\_id. Report unmatched rows.'' \\[2pt]
    Validate &
      ``Check for duplicate (state, year) pairs and summarize missing values by group.'' \\[2pt]
    Plot &
      ``Plot mean employment by treatment group over time, vertical line at policy date.'' \\
    \bottomrule
  \end{tabular}
\end{frame}

\begin{frame}{What Claude Code Gets Wrong}
  \small
  Claude Code is helpful but not infallible. Common failure modes:

  \vspace{2pt}

  \begin{enumerate}
    \item \textbf{Hallucinated variable names} -- \textit{Fix:} have it read the data first
    \item \textbf{Silent data loss} -- a merge may drop rows. \textit{Fix:} compare row counts before/after
    \item \textbf{Wrong join type} -- inner when you wanted left. \textit{Fix:} specify explicitly; use \texttt{anti\_join()}
    \item \textbf{Outdated syntax} -- deprecated functions. \textit{Fix:} run the code and check warnings
  \end{enumerate}

  \vspace{4pt}

  \normalsize
  \textbf{\textcolor{usfgreen}{Rule:}} always review the code Claude Code writes. You are the researcher; it is the assistant.

  \vspace{2pt}

  \textit{Have you ever had a merge silently drop observations? How would you detect that?}
\end{frame}

\begin{frame}{Managing Claude Code's Memory}
  \small
  Claude Code has \textbf{limited context} -- it forgets between sessions and can lose track within a long one.

  \vspace{3pt}

  \textbf{How to work around this:}
  \begin{enumerate}\setlength{\itemsep}{1pt}
    \item \textbf{Commit often with clear messages.} Git is Claude Code's long-term memory -- it can read commits and diffs to recover context.
    \item \textbf{Keep CLAUDE.md current.} Loaded every session; put anything important here.
    \item \textbf{Use \texttt{.claude/rules/}} for conventions you never want to re-explain.
    \item \textbf{Paste context} when starting fresh: ``The script is \texttt{code/01\_clean.R}. Pick up where we left off.''
    \item \textbf{Tell it to read files} rather than summarize from memory.
  \end{enumerate}

  \vspace{2pt}

  \textbf{\textcolor{usfgreen}{Rule:}} if Claude Code seems confused, it lost context. Point it to specific files and commits.
\end{frame}

% ══════════════════════════════════════════════════════════
%  SECTION 5: Work Sprint
% ══════════════════════════════════════════════════════════

\section{Work Sprint}

\begin{frame}{Work Sprint Goals}
  \textbf{\textcolor{usfgreen}{35 minutes.}} Work on your own project data with Claude Code.

  \vspace{6pt}

  Choose the task that matches where you are:

  \vspace{4pt}

  \small
  \begin{tabular}{@{}l>{\raggedright\arraybackslash}p{8.5cm}@{}}
    \toprule
    \textbf{If you have...} & \textbf{Do this:} \\
    \midrule
    Raw data loaded &
      Clean it: handle missing values, standardize codes, reshape to panel format \\[4pt]
    Multiple datasets &
      Merge them: identify keys, run joins, validate match rates \\[4pt]
    Clean data already &
      Validate: run the five checks from the validation slide, then start exploratory plots \\[4pt]
    No data yet &
      Use Claude Code to help you find and download a relevant public dataset \\
    \bottomrule
  \end{tabular}

  \normalsize
  \vspace{6pt}

  \textbf{Goal:} by the end of this sprint, you should have a working cleaning script committed to Git.
\end{frame}

\begin{frame}{Sprint Checkpoint}
  \textbf{Before you stop:}

  \vspace{8pt}

  \begin{enumerate}
    \item Did you write a \textbf{cleaning script} (not just run ad hoc commands)?
    \item Can you delete \texttt{data/cleaned/} and \textbf{regenerate} it from the script?
    \item Did you \textbf{validate} your cleaned data (row count, no duplicates, plausible values)?
    \item Did you \textbf{commit and push} to Git?
    \item Does your \textbf{CLAUDE.md} describe the data pipeline?
  \end{enumerate}

  \vspace{8pt}

  If you answered ``no'' to any of these, take 5 more minutes to fix it.
\end{frame}

% ══════════════════════════════════════════════════════════
%  Share-outs & Wrap-up
% ══════════════════════════════════════════════════════════

\begin{frame}{Share-outs}
  \textbf{3--4 students, 3 minutes each.}

  \vspace{12pt}

  Cover:
  \begin{enumerate}
    \item \textbf{What data problems did you find?}
    \begin{itemize}
      \item Missing values, inconsistent coding, merge failures, etc.
    \end{itemize}
    \item \textbf{How did Claude Code help?}
    \begin{itemize}
      \item What did you ask it to do? Did it work on the first try?
    \end{itemize}
    \item \textbf{What is your data pipeline now?}
    \begin{itemize}
      \item Raw $\to$ cleaning script $\to$ cleaned data $\to$ ready for analysis?
    \end{itemize}
  \end{enumerate}

  \vspace{12pt}

  Listeners: one tip or one question for each presenter.
\end{frame}

\begin{frame}{Key Takeaways}
  \begin{enumerate}
    \item \textbf{CLAUDE.md} is the single most important file for using Claude Code effectively -- it gives context that persists across sessions

    \vspace{4pt}

    \item \textbf{Raw data is sacred}: every transformation must be in a script, never done by hand

    \vspace{4pt}

    \item \textbf{Validate after every step}: row counts, duplicates, merge diagnostics, range checks

    \vspace{4pt}

    \item \textbf{Claude Code accelerates the tedious parts} -- but you are still responsible for the research design decisions
  \end{enumerate}
\end{frame}

\begin{frame}{Wrap-Up \& Next Week}
  \textbf{Due this week:}
  \begin{itemize}
    \item Weekly progress report (Wednesday)
  \end{itemize}

  \vspace{8pt}

  \textbf{Next week (Week 6):}
  \begin{itemize}
    \item Exploratory data analysis and early figures
    \item \textbf{Empirical Strategy Draft due Friday, March 6}
  \end{itemize}

  \vspace{8pt}

  \textbf{Before next class:}
  \begin{itemize}
    \item Finish your cleaning/merging pipeline and push scripts + CLAUDE.md to GitHub
    \item Start exploring: summary stats and one or two preliminary figures
    \item Begin drafting your empirical strategy section (builds on Week 4 DAG work)
  \end{itemize}
\end{frame}

\end{document}
